\subsection{Listings -- Source Code}

The \texttt{listings} package is used to display source code while
preserving indentation and formatting. Different styles can be used
depending on the programming language.

For example, MATLAB code can be displayed using the
\texttt{matlab} style:

\begin{lstlisting}[style=matlab]
clc
clear

x = linspace(0, 2*pi, 100);
y = sin(x);

plot(x,y)
\end{lstlisting}

The corresponding environment is opened with:

\begin{lstlisting}[style=latex]
\begin{lstlisting}[style=matlab]
\end{lstlisting}

and closed using \verb|\end{lstlisting}|.

The appearance of the code listings is configured in
\texttt{Packages/Code/package\_config.tex}. This file defines the
available styles, including properties such as the font, colors,
background, line numbering and spacing.

When adding or modifying a code style, it is recommended to make the
changes in this configuration file rather than changing individual
listings throughout the document.


